
\documentclass{article}
\usepackage{colortbl}
\usepackage{makecell}
\usepackage{multirow}
\usepackage{supertabular}

\begin{document}

\newcounter{utterance}

\centering \large Interaction Transcript for game `eqbench', experiment `eq\_bench', episode 21 with reflexion{-}llama{-}3.1{-}8b{-}instruct.
\vspace{24pt}

{ \footnotesize  \setcounter{utterance}{1}
\setlength{\tabcolsep}{0pt}
\begin{supertabular}{c@{$\;$}|p{.15\linewidth}@{}p{.15\linewidth}p{.15\linewidth}p{.15\linewidth}p{.15\linewidth}p{.15\linewidth}}
    \# & \multicolumn{2}{c}{Player} && \multicolumn{2}{c}{Game Master} \\
    \hline

    \theutterance \stepcounter{utterance}  
    & & & \multicolumn{4}{p{0.6\linewidth}}{
        \cellcolor[rgb]{0.9,0.9,0.9}{
            \makecell[{{p{\linewidth}}}]{
                \texttt{\tiny{[P1$\langle$GM]}}
                \texttt{You have attempted this same task before and did not succeed. These are the plans you wrote after those attempts.} \\
\texttt{Trial \#1: New plan: I overestimated the intensity of Elliot's emotions, particularly the resigned and embarrassed scores. Upon re{-}reading the dialogue, I realize that Elliot's tone is more hesitant and uncertain, which suggests that he is more likely to feel embarrassed and less likely to feel resigned. Elliot's words "I have feelings for you" and "more than just friendship" indicate a sense of vulnerability and nervousness, which aligns more with the feeling of embarrassment. In contrast, Elliot's tone does not convey a sense of resignation or hopelessness. Elliot's words also do not suggest that he is angry, as he is being open and honest about his feelings. Therefore, I will adjust my scores to reflect Elliot's more vulnerable and uncertain tone.} \\
\\ 
\\ 
\\ 
\texttt{Your task is to predict the likely emotional responses of a character in this dialogue:} \\
\\ 
\texttt{Elliot: Alex, there's something I need to tell you.} \\
\texttt{Alex: Sure, man. What's up?} \\
\texttt{Elliot: I... I've been struggling with something. Something I can't keep to myself anymore.} \\
\texttt{Alex: Okay, you're starting to worry me. What is it?} \\
\texttt{Elliot: It's about... me. And my feelings. My feelings for... someone.} \\
\texttt{Alex: Well, who is it? Sophie? She's pretty cute.} \\
\texttt{Elliot: No, it's not Sophie. It's... It's you, Alex.} \\
\texttt{Alex: Me? What do you mean?} \\
\texttt{Elliot: I mean I have feelings for you. More than just friendship.} \\
\texttt{Alex: But... I'm with Sophie.} \\
\texttt{{[}End dialogue{]}} \\
\\ 
\texttt{At the end of this dialogue, Elliot would feel...} \\
\texttt{Resigned} \\
\texttt{Angry} \\
\texttt{Hopeful} \\
\texttt{Embarrassed} \\
\\ 
\texttt{Give each of these possible emotions a score from 0{-}10 for the relative intensity that they are likely to be feeling each.} \\
\\ 
\texttt{You must output in the following format, including headings (of course, you should give your own scores), with no additional commentary:} \\
\\ 
\texttt{Resigned: <score>} \\
\texttt{Angry: <score>} \\
\texttt{Hopeful: <score>} \\
\texttt{Embarrassed: <score>} \\
\\ 
\\ 
\texttt{{[}End of answer{]}} \\
\\ 
\texttt{Remember: zero is a valid score, meaning they are likely not feeling that emotion. You must score at least one emotion > 0.} \\
\\ 
\texttt{Your answer:} \\
            }
        }
    }
     \\ \\

    \theutterance \stepcounter{utterance}  
    & \multicolumn{4}{p{0.6\linewidth}}{
        \cellcolor[rgb]{0.561,0.737,0.561}{
            \makecell[{{p{\linewidth}}}]{
                \texttt{\tiny{[P1$\rangle$GM]}}
                \texttt{Resigned: 0} \\
\texttt{Angry: 0} \\
\texttt{Hopeful: 0} \\
\texttt{Embarrassed: 8} \\
            }
        }
    }
    & & \\ \\

    \theutterance \stepcounter{utterance}  
    & & & \multicolumn{2}{p{0.3\linewidth}}{
        \cellcolor[rgb]{0.9,0.9,0.9}{
            \makecell[{{p{\linewidth}}}]{
                \texttt{\tiny{[GM$|$GM]}}
                \texttt{Resigned: 0} \\
\texttt{Angry: 0} \\
\texttt{Hopeful: 0} \\
\texttt{Embarrassed: 8} \\
            }
        }
    }
    & & \\ \\

    \theutterance \stepcounter{utterance}  
    & & & \multicolumn{2}{p{0.3\linewidth}}{
        \cellcolor[rgb]{0.9,0.9,0.9}{
            \makecell[{{p{\linewidth}}}]{
                \texttt{\tiny{[GM$|$GM]}}
                \texttt{Resigned: 7} \\
\texttt{Angry: 0} \\
\texttt{Hopeful: 3} \\
\texttt{Embarrassed: 8} \\
            }
        }
    }
    & & \\ \\

    \theutterance \stepcounter{utterance}  
    & & & \multicolumn{2}{p{0.3\linewidth}}{
        \cellcolor[rgb]{0.9,0.9,0.9}{
            \makecell[{{p{\linewidth}}}]{
                \texttt{\tiny{[GM$|$GM]}}
                \texttt{game\_result = LOSE} \\
            }
        }
    }
    & & \\ \\

    \theutterance \stepcounter{utterance}  
    & & \multicolumn{4}{p{0.6\linewidth}}{
        \cellcolor[rgb]{0.95,0.95,0.95}{
            \makecell[{{p{\linewidth}}}]{
                \texttt{\tiny{[P1$\leftrightarrow$P]}}
                \texttt{Post{-}episode Reflexion (not sent to GM):} \\
\texttt{New plan: I overestimated the intensity of Elliot's emotions, particularly the resigned and hopeful scores. Upon re{-}reading the dialogue, I realize that Elliot's tone is more hesitant and uncertain, which suggests that he is more likely to feel embarrassed and less likely to feel resigned or hopeful. Elliot's words "I have feelings for you" and "more than just friendship" indicate a sense of vulnerability and nervousness, which aligns more with the feeling of embarrassment. In contrast, Elliot's tone does not convey a sense of resignation or hopelessness. Elliot's words also do not suggest that he is angry, as he is being open and honest about his feelings. However, I also realize that Elliot might feel a sense of disappointment or sadness due to Alex's current relationship with Sophie. Therefore, I will adjust my scores to reflect Elliot's more vulnerable and uncertain tone, and also consider the potential disappointment or sadness he might feel. } \\
\\ 
\texttt{Resigned: 2} \\
\texttt{Angry: 0} \\
\texttt{Hopeful: 1} \\
\texttt{Embarrassed: 9} \\
            }
        }
    }
    & \\ \\

\end{supertabular}
}

\end{document}
